\begin{abstract}
The classical polynomial
\((X^3-19)(X^2+X+1)\) has a root over every completion of \(\Q\), but no
rational root, and at first looks like an isolated arithmetic curiosity.
Starting from this Hasse failure, we ask whether
the map, rather than only the fiber, can be kept fixed.  The decisive
observation is that after centering the quadratic factor, the whole family
\((X^3-a)(X^2+X+1)\) moves in exactly the two coefficients controlled by
one quadratic-gauge target line.

This produces one explicit polynomial map
\(\Phi:\A^3_{\Q}\to\A^3_{\Q}\) with Jacobian determinant one and geometric
degree five.  For every noncube integer \(a>1\) such that
\[
 a\equiv1\pmod 9,
 \qquad p\mid a\Longrightarrow p\equiv1\pmod3,
\]
the map has a full fiber
\[
 \Spec\Q[X]/\bigl((X^3-a)(X^2+X+1)\bigr).
\]
Every such fiber is reduced, has points over \(\R\) and every \(\Q_p\), and
has no \(\Q\)-point.  The targets lie on a fixed rational line and have
height \(32a\).  The local argument does not require \(a\) to be prime:
its congruence-and-support core is multiplicatively closed, after which
the perfect cubes are removed.  A character-twisted
Selberg--Delange argument then shows that the number of constructed targets
of height at most \(B\) is asymptotic to
\[
 \frac{G_3(1)}{96\sqrt{\pi}}\,
 \frac{B}{\sqrt{\log B}},
\]
where \(G_3(1)>0\) is an explicit Euler product.  Finally, the classical
degree bound for intersective polynomials shows that geometric degree five
is best possible.  Thus one exceptional quintic becomes a quantitatively
large family of failures on a single map of the smallest possible geometric
degree.
\end{abstract}
